\documentclass[11pt]{article}
\usepackage{amsmath,amssymb,amsthm}
\usepackage[margin=1in]{geometry}

\newtheorem{theorem}{Theorem}
\newtheorem{lemma}[theorem]{Lemma}
\newtheorem{proposition}[theorem]{Proposition}
\newtheorem{corollary}[theorem]{Corollary}
\newtheorem{conjecture}[theorem]{Conjecture}
\newtheorem{remark}[theorem]{Remark}

\DeclareMathOperator{\Corr}{Corr}
\DeclareMathOperator{\Cov}{Cov}
\DeclareMathOperator{\Var}{Var}

\title{Why the Peak--Gap Correlation is Positive:\\A Proof Sketch}
\author{Working notes}
\date{March 2026}

\begin{document}
\maketitle

\section{The Claim}

\begin{conjecture}\label{conj:main}
For any $N \geq N_0$, the trigonometric sum
\[
f(t) = \sum_{n=1}^{N} \frac{\cos(\omega_n t + \phi_n)}{\sqrt{n}},
\qquad \omega_n = \log(n+1),
\]
has peak--gap correlation $r = \Corr(g_k, P_k) > 0$, uniformly in the phases $\phi_n \in [0, 2\pi)$.
Moreover, for random i.i.d.\ uniform phases, $\mathbb{E}[r] \to r_\infty > 0$ as $N \to \infty$.
\end{conjecture}

\noindent
Numerical evidence: $\mathbb{E}[r] \approx 0.13$ for $N = 20$ to $1000$.

\section{Two-Layer Structure}

The correlation $r$ arises from two layers:

\subsection{Layer 1: The MVT mechanism (universal)}
Between consecutive zeros $\gamma_k$ and $\gamma_{k+1}$, separated by gap $g_k$:
\[
P_k = \max_{t \in [\gamma_k, \gamma_{k+1}]} |f(t)|.
\]
By the mean-value theorem and Taylor expansion at $\gamma_k$:
\[
P_k \approx \frac{|f'(\gamma_k)| \cdot g_k}{c_k},
\]
where $c_k \in [\pi, 2]$ is a shape factor depending on the curvature.

If $X_k = |f'(\gamma_k)|/c_k$ is independent of $g_k$:
\[
\Cov(g_k, P_k) = \Cov(g_k, X_k g_k) = \mathbb{E}[X_k] \cdot \Var(g_k) > 0.
\]

This gives a positive lower bound on $r$ whenever:
\begin{enumerate}
\item $\Var(g) > 0$ (gaps are not all equal);
\item $\mathbb{E}[|f'(\gamma)|] > 0$ (trivially true);
\item $\Corr(|f'(\gamma)|, g)$ is not too negative.
\end{enumerate}

\textbf{Verified}: For RS-like sums with random phases, $\Corr(|f'|, g) \approx 0.00 \pm 0.01$.
The near-independence of $|f'|$ and $g$ is a consequence of the high-dimensional phase space:
with many incommensurate frequencies, the derivative at a zero is effectively independent
of the local gap structure.

\subsection{Layer 2: The saturation correction (reduces $r$)}

For large gaps ($g > 1.5\bar{g}$), the MVT approximation breaks down:
the function oscillates multiple times within the gap, and the peak saturates.
This creates a hump-shaped $\mathbb{E}[P \mid g]$, reducing $r$ from the MVT prediction
of $\approx 0.71$ to the actual $\approx 0.13$.

Crucially, the saturation only affects the \emph{tail} of the gap distribution,
while the positive MVT contribution comes from the \emph{bulk}.
Since the bulk dominates (gaps $> 1.5\bar{g}$ are rare), $r$ remains positive.

\section{Why $r > 0$ for RS-Type Sums but Not All Sums}

Counter-example: $f(t) = \cos(t) + \cos(2t)$ has $r = -1.0$.
Here the frequencies are commensurate (ratio 2:1), creating a periodic function
where the gap-peak relationship is inverted.

For RS-like sums, the frequencies $\omega_n = \log(n+1)$ are highly incommensurate:
$\log(n+1)/\log(m+1)$ is irrational for $n \neq m$ (by the fundamental theorem of arithmetic,
since this would require $(n+1)^a = (m+1)^b$ for integer $a, b$).

This incommensurability ensures:
\begin{enumerate}
\item The process $f(t)$ is ergodic (almost periodic, not periodic);
\item The zero locations are ``generic'' (not organized by a lattice);
\item The derivative $|f'(\gamma)|$ is effectively independent of the gap $g$.
\end{enumerate}

\section{Proof Strategy}

\begin{enumerate}
\item \textbf{CLT for $f(t)$}: For large $N$, the sum $f(t) = \sum a_n \cos(\omega_n t + \phi_n)$
converges to a Gaussian process with correlation function
$C(\tau) = \sum (1/n) \cos(\omega_n \tau) / \sum (1/n)$.

\item \textbf{Spectral moments}: The moments $m_k = \sum (\omega_n)^k / n$ are:
\[
m_0 = \sum \frac{1}{n} \approx \log N, \quad
m_2 = \sum \frac{(\log n)^2}{n} \approx \frac{(\log N)^3}{3}, \quad
m_4 = \sum \frac{(\log n)^4}{n} \approx \frac{(\log N)^5}{5}.
\]

\item \textbf{Rice's formula}: The mean zero density is $\lambda_0 = \sqrt{m_2/m_0}/\pi$
and the mean extremum density is $\lambda_2 = \sqrt{m_4/m_2}/\pi$.

\item \textbf{Excursion theory}: For a Gaussian process with these spectral moments,
$\mathbb{E}[P \mid g]$ is computable from the conditional distribution
of the process given zeros at the endpoints.
The key result (from Cram\'er--Leadbetter theory):
\[
\mathbb{E}[P \mid g] = \sigma \cdot h(g/\ell),
\]
where $\sigma = \sqrt{m_0}$ is the RMS, $\ell = \pi\sqrt{m_0/m_2}$ is the correlation length,
and $h$ is a universal function that increases for $g < g^*$ and saturates for $g > g^*$.

\item \textbf{Positive covariance}: Since $h$ is increasing for the bulk of the gap distribution
and the saturation only affects the tail:
\[
\Cov(g, P) = \sigma \cdot \Cov(g, h(g/\ell)) > 0.
\]

\item \textbf{Uniformity in $N$}: As $N \to \infty$, both $\sigma$ and $\ell$ change, but
the \emph{normalized} quantities $g/\ell$ and $P/\sigma$ have limiting distributions.
The function $h$ depends only on the ratio $m_0 m_4 / m_2^2$ (the roughness parameter),
which converges to $5/3$ for the RS spectral density.
Therefore $r = \Corr(g, P)$ converges to a positive limit $r_\infty > 0$.
\end{enumerate}

\section{From Random Phases to $\zeta$}

For the actual Riemann zeta function, the phases $\phi_n = \theta(t) - t\log n$ are
\emph{not} random---they are determined by the functional equation.
These specific phases give $r \approx 0.63$--$0.93$, far above the random-phase
prediction of $r \approx 0.13$.

The enhancement comes from the functional equation phase $\theta(t)$, which creates
a positive correlation $\Corr(|\zeta'|, G) \approx 0.81$ between the derivative at zeros
and the gap (the ``self-correcting mechanism'').
This is absent for random phases, where $\Corr(|f'|, g) \approx 0$.

For the density bound, we only need $r > 0$, which holds for both random-phase
and specific-phase RS sums.
The stronger bound from the specific phase is a bonus.

\section{Open Problems}

\begin{enumerate}
\item Make step 5 rigorous: prove $\Cov(g, h(g/\ell)) > 0$ for the specific function $h$
arising from the RS spectral density.
\item Prove the CLT convergence (step 1) with quantitative error bounds in the $r$ statistic.
\item Show the roughness parameter $m_0 m_4/m_2^2$ is bounded, ensuring the limit exists.
\end{enumerate}

\end{document}
